\section{Conclusion}
\label{sec:conclusion}

Algorithmically equivalent student submissions using different programming languages constitute an interesting type of structured similarity that CPG-based normalization is ideally suited to identify. By addressing the CES labeling incompatibility issue that limited cross-frontend comparisons in previously developed language-specific CPG models, the proposed Context Normalization Protocol allows direct comparison of patterns among C, C++, and Java, without learning any alignments---resulting in interpretable, submission-level algorithmic inventories that are readily usable to facilitate feedback in resource-constrained educational settings.

Empirically, we find that the aggregate cross-language retrieval accuracy is 87.25\% ($349/400$), with the core directions achieving 88.60\% ($342/386$) and the Java$\rightarrow$C/C++ direction exhibiting 91.43\% ($128/140$) top-1 recall. This aggregate performance is stratified by directional asymmetry (detailed in Section~\ref{subsec:cross_lang_results}), which serves as the primary boundary condition of our empirical results. The feature weights, trained only using monolingual C samples, prove invariant throughout our mixed-language evaluations. When compared against a fine-tuned UniXcoder model in a setting that ensures an apples-to-apples comparison, the deterministic method exhibits a 17.38~percentage point (pp) improvement over the UniXcoder model for this low-resource setting but not necessarily under any arbitrary condition. In particular, compared against a standalone BL-view ablation, the fused framework provides a statistically significant 7.75~pp benefit (McNemar $p{<}0.0001$, OR$=4.875$ [95\% CI: 2.26, 10.52]), indicating the added value of fusing structural and semantic perspectives with lexical features.
The following four boundary conditions have been specified for the applicability scope of the above results: firstly, the directionality bias of the parser front-end described earlier; secondly, the restricted generalizability due to the use of data from a single institution; thirdly, the relatively low reference count per problem within the cross-language reference corpus (6–7 references per query); and fourthly, the limited scope of application of the CES taxonomy applicable only to CS-1 imperative constructs. The APM generator is merely a proof-of-feasibility under author-constructed test conditions; the agreement percentage of 98.3\% should therefore be seen strictly as a consistency measure and not an established performance benchmark. All these limitations are discussed further in Section~\ref{subsec:threats}.

The gradient-free nature of the method makes this framework applicable as an educational analysis tool that does not require either GPU hardware or corpora for fine-tuned learning. Future research should focus on multi-institution validation, characterizing the growth behavior of the reference pool, and testing the capabilities of the APM using third-party evaluations.

\subsection{Future Work}
\label{subsec:future_work}

Future directions include:
\begin{enumerate}
    \item  Multi-institution validation of the similarity retrieval framework using external sources for ground truth, specifically collecting balanced C$\rightarrow$Java and C++$\rightarrow$Java query populations across institutions so that the underpowered directions can be evaluated with inferential validity.
    \item Extending the Context Normalization Protocol and CES taxonomy to dynamic languages (Python, JavaScript), where the CPG representation deviates even further from that of C-like language code.
    \item Extending the CES taxonomy to include current idioms in modern C++17/20 (range pipelines, structured bindings, \texttt{std::ranges::views}) and Java stream collectors.
    \item Using fuzz tests from third parties on the rule generation engine using an extended CS-1 problem corpus and external test suites for determining failure modes.
    \item Cross-validation of the Java and cross-language samples separately to determine multi-lingual weight stability without relying on the monolingual CV results.
    \item Testing tuned neural baselines against a larger labeled data set to empirically determine the limitations of learning approaches compared to deterministic methods as training data is increased.
    \item Benchmarking against CPG code assessment frameworks and existing cross-language retrieval research as the literature in this field expands.
\end{enumerate}

\subsection{Open Science Statement}

The dataset, CPG generation script, APM generator modules, the evaluation framework, weight grid search settings, and the Joern version and JVM settings used during evaluation are all publicly available at:
\url{https://github.com/Abhinav-C-Das/Multi-language-extension-of-semantic-code-similarity-framework-with-code-translation}

The source code includes the CES extraction scripts, Java API mapping information, STL pattern mapping table, APM generator modules, structured test suites, and the 23-pattern entry point list.

\section*{Acknowledgment}
The authors gratefully acknowledge Sri Mata Amritanandamayi Devi (Amma), Chancellor, Amrita Vishwa Vidyapeetham, for her inspiration and for providing financial support for the Article Processing Charges (APC) of this publication.